\section{Clean modelling dataset and train-test split}

\subsection{Purpose of this stage}

The raw-data audit produced an interim dataset in which the representation problem in \texttt{TotalCharges} was corrected. The next step is to create the supervised modelling table and the held-out train-test split.

This stage is deliberately narrow. It creates the binary target, defines the feature set, validates semantic feature groups, and creates a stratified split. It does not perform exploratory feature-target analysis, one-hot encoding, scaling, feature engineering, resampling, feature selection, or model fitting. Those steps are postponed until after the split and should use the training set only.

\subsection{Interim dataset}

The interim dataset is loaded from:

\[
\texttt{data/interim/telco\_churn\_interim.csv}.
\]

It has the same raw columns as the original dataset, except that \texttt{TotalCharges} has been converted to a numeric representation and the zero-tenure blank values have been set to \texttt{0.0}. No remaining missing values are present at this stage.

\begin{table}[H]
\centering
\caption{Interim dataset overview}
\begin{tabular}{lr}
\toprule
Item & Value \\
\midrule
Number of rows & 7043 \\
Number of columns & 21 \\
Duplicate rows & 0 \\
Total missing values & 0 \\
Columns with missing values & 0 \\
\bottomrule
\end{tabular}
\end{table}

\subsection{Binary target encoding}

The supervised learning target is \texttt{Churn\_binary}. The positive class is churn, because this is the event of practical interest.

\[
\texttt{Churn\_binary} =
\begin{cases}
1, & \text{if } \texttt{Churn = Yes},\\
0, & \text{if } \texttt{Churn = No}.
\end{cases}
\]

The encoding produces no missing values and preserves the two-class structure of the original target.

\begin{table}[H]
\centering
\caption{Target encoding}
\begin{tabular}{lr}
\toprule
Original target level & Encoded value \\
\midrule
No & 0 \\
Yes & 1 \\
\bottomrule
\end{tabular}
\end{table}

\subsection{Modelling feature set}

The modelling feature set excludes \texttt{customerID}, the original string target \texttt{Churn}, and the binary target \texttt{Churn\_binary}. The identifier is excluded because it is row-specific and does not represent a generalizable customer characteristic. The original target is excluded because the model uses the encoded binary target.

The resulting modelling table contains 7043 observations, 19 feature columns, and one target column.

\begin{table}[H]
\centering
\caption{Clean modelling table overview}
\begin{tabular}{lr}
\toprule
Item & Value \\
\midrule
Rows & 7043 \\
Feature columns & 19 \\
Target columns & 1 \\
Identifier included as feature & No \\
Original string target included as feature & No \\
\bottomrule
\end{tabular}
\end{table}

\subsection{Semantic feature groups}

The 19 modelling features are divided into semantic groups. These groups describe the clean base dataset; they do not yet define the final preprocessing for every model. Model-specific preprocessing is handled later inside training-only pipelines.

\begin{table}[H]
\centering
\caption{Feature groups for modelling}
\begin{tabular}{lrp{9.5cm}}
\toprule
Feature group & Count & Features \\
\midrule
Numeric features & 3 & \texttt{tenure}, \texttt{MonthlyCharges}, \texttt{TotalCharges} \\
Binary categorical features & 6 & \texttt{SeniorCitizen}, \texttt{gender}, \texttt{Partner}, \texttt{Dependents}, \texttt{PhoneService}, \texttt{PaperlessBilling} \\
Nominal categorical features & 10 & \texttt{MultipleLines}, \texttt{InternetService}, \texttt{OnlineSecurity}, \texttt{OnlineBackup}, \texttt{DeviceProtection}, \texttt{TechSupport}, \texttt{StreamingTV}, \texttt{StreamingMovies}, \texttt{Contract}, \texttt{PaymentMethod} \\
\bottomrule
\end{tabular}
\end{table}

The feature-group validation confirms that every modelling feature belongs to exactly one group: no feature is missing from the groups, no grouped feature lies outside the modelling feature set, and no feature is assigned more than once.

\subsection{Final checks before splitting}

The final modelling table contains no missing values. Exact duplicate rows in the full interim data are absent. After excluding the identifier, duplicate feature-target combinations may exist, but this is not automatically a data-quality issue because multiple customers can genuinely share the same observed characteristics and churn label.

\begin{table}[H]
\centering
\caption{Final checks before splitting}
\begin{tabular}{lr}
\toprule
Item & Value \\
\midrule
Missing values in modelling table & 0 \\
Duplicate full interim rows & 0 \\
Duplicate feature-target rows after excluding identifier & 22 \\
\bottomrule
\end{tabular}
\end{table}

\subsection{Target distribution before splitting}

The target distribution is moderately imbalanced. The positive class, churn, represents approximately 26.5\% of observations. This motivates a stratified train-test split. This target-only summary is used for split design rather than feature selection or model interpretation. No feature-target relationships are inspected before the train-test split.

\begin{table}[H]
\centering
\caption{Target distribution before splitting}
\begin{tabular}{lrr}
\toprule
\texttt{Churn\_binary} & Count & Percentage \\
\midrule
0 & 5174 & 73.46 \\
1 & 1869 & 26.54 \\
\bottomrule
\end{tabular}
\end{table}

\subsection{Train-test split strategy}

The clean modelling table is split into a training set and a held-out test set. The test set is reserved for final evaluation and should not be used for feature-target EDA, preprocessing design, model selection, hyperparameter tuning, or threshold tuning.

The split uses:

\[
\texttt{test\_size} = 0.20,
\qquad
\texttt{random\_state} = 42.
\]

The split is stratified by \texttt{Churn\_binary}, so the train and test sets preserve the original churn proportion.

\begin{table}[H]
\centering
\caption{Train-test split overview}
\begin{tabular}{lrr}
\toprule
Split & Rows & Percentage \\
\midrule
Train & 5634 & 79.99 \\
Test & 1409 & 20.01 \\
\bottomrule
\end{tabular}
\end{table}

\subsection{Target distribution by split}

The stratified split preserves the target distribution closely.

\begin{table}[H]
\centering
\caption{Target distribution by split}
\begin{tabular}{llrr}
\toprule
Split & \texttt{Churn\_binary} & Count & Percentage \\
\midrule
Train & 0 & 4139 & 73.46 \\
Train & 1 & 1495 & 26.54 \\
Test & 0 & 1035 & 73.46 \\
Test & 1 & 374 & 26.54 \\
\bottomrule
\end{tabular}
\end{table}

\subsection{Saved outputs}

The training and test datasets are saved locally as:

\[
\texttt{data/processed/train.csv},
\qquad
\texttt{data/processed/test.csv}.
\]

These files form the boundary for the rest of the project. The next workflow stage should load only \texttt{train.csv} for exploratory analysis and modelling decisions. The held-out test file should remain unused until final model evaluation.
